\documentclass[11pt]{article}

\usepackage[margin=1in]{geometry}
\usepackage{amsmath,amssymb,amsthm,mathtools}
\usepackage{enumitem}
\usepackage{microtype}
\usepackage{hyperref}

\newtheorem{conjecture}{Conjecture}

\title{Proposed Direction: Power-Saving Compression for Nonsingular Quadratic Polynomials}
\author{}
\date{}

\begin{document}

\maketitle

The recent experiments suggest a natural next question in the Elekes--R\'onyai direction.

The known 2026 counterexamples proving a fixed power saving
\[
|P(A,A)|\le |A|^{2-c}
\]
use polynomials whose quadratic part is \emph{rank one}, essentially of the form
\[
(x-y)^2+\text{lower-order terms}.
\]

For example, the mechanism for
\[
P(x,y)=5(x+y)+(x-y)^2
\]
works because, after introducing a suitable denominator \(Q\),
\[
Q^2P(b/Q,c/Q)
=
(b-c)^2+5Q(b+c),
\]
so modulo every prime \(q\mid Q\) the lower-order part disappears and the output is forced to be a quadratic residue:
\[
Q^2P(b/Q,c/Q)\equiv (b-c)^2\pmod q.
\]

Completely split primes then replicate this restriction independently across every residue coordinate of a high-degree number field. The Chinese remainder theorem gives an exponential saving in the degree, while bounded root discriminant keeps the lattice-counting loss only exponential. This is closely analogous to the high-degree amplification mechanism in the recent unit-distance construction.

However, this mechanism fundamentally exploits the degeneracy
\[
\det
\begin{pmatrix}
2&-2\\
-2&2
\end{pmatrix}
=0.
\]

This suggests asking whether fixed power-saving compression persists once the quadratic part is genuinely two-dimensional.

\section{Conjecture 1: Canonical Test Case}

Let
\[
Q(x,y)=x^2+xy+y^2.
\]

\begin{conjecture}[Eisenstein Norm Compression]
There exists an absolute constant \(c>0\) and arbitrarily large finite sets
\[
A\subset\mathbb R
\]
such that
\[
\boxed{
|Q(A,A)|\le |A|^{2-c}.
}
\]
\end{conjecture}

This is the immediate conjecture proposed for investigation.

It is structurally different from the known counterexamples because
\[
\nabla^2Q=
\begin{pmatrix}
2&1\\
1&2
\end{pmatrix},
\qquad
\det\nabla^2Q=3\neq0.
\]

Thus the quadratic part has full rank, and the existing denominator trick cannot make it collapse to a one-dimensional square map.

There is also a natural number-theoretic structure. If
\[
\omega^2+\omega+1=0,
\]
then
\[
x^2+xy+y^2
=
N_{\mathbb Q(\omega)/\mathbb Q}(x-y\omega).
\]

So this is the Eisenstein-norm analogue of the Gaussian norm \(x^2+y^2\), making it a particularly natural polynomial on which to try to transfer ideas from the unit-distance construction.

The classical set \(A=[n]\) already gives logarithmic compression of this form, roughly
\[
|Q(A,A)|\ll \frac{n^2}{\sqrt{\log n}},
\]
but the proposed conjecture asks for the qualitatively stronger bound
\[
|Q(A,A)|\le n^{2-c}
\]
with fixed \(c>0\).

\section{Why This Is Not Covered by the Current CRT Mechanism}

The previous experiments also give a useful obstruction.

For a homogeneous quadratic such as
\[
Q(x,y)=x^2+xy+y^2,
\]
rescaling does not create a new degeneration:
\[
Q(b/Q,c/Q)=\frac{Q(b,c)}{Q^2}.
\]

After clearing denominators, one recovers exactly the same quadratic form:
\[
Q^2Q(b/Q,c/Q)=Q(b,c).
\]

Thus there are no lower-degree terms that disappear modulo primes dividing \(Q\).

Moreover, the ordinary local Cartesian restriction experiment gives inequalities of the type
\[
R^k\,|Q(B,B)|\ge |B|^2,
\]
so simply restricting the input modulo split prime powers cannot produce normalized compression.

Therefore a proof of the conjecture would require a genuinely different use of the arithmetic structure rather than another application of the rank-one square-residue sieve.

\section{Conjecture Hierarchy}

There are three natural levels.

\subsection{Level 1: Eisenstein Norm Compression}

Prove the concrete statement
\[
\boxed{
|x^2+xy+y^2(A,A)|\le |A|^{2-c}
}
\]
for arbitrarily large \(A\subset\mathbb R\).

This is the first target.

\subsection{Level 2: Binary Norm-Form Compression}

Let
\[
q(x,y)=ax^2+bxy+cy^2
\]
be a positive-definite binary quadratic form over \(\mathbb Q\), with
\[
b^2-4ac<0,
\]
and assume that \(q\) is not of an Elekes--R\'onyai additive or multiplicative special form.

\begin{conjecture}[Binary Norm-Form Compression]
There exists \(c_q>0\) and arbitrarily large finite sets
\[
A\subset\mathbb R
\]
such that
\[
\boxed{
|q(A,A)|\le |A|^{2-c_q}.
}
\]
\end{conjecture}

This asks whether the phenomenon is associated with quadratic norm forms rather than only the single Eisenstein example.

\subsection{Level 3: Nonsingular Quadratic Compression}

Let
\[
f(x,y)=ax^2+bxy+cy^2+dx+ey+r\in\mathbb Q[x,y]
\]
be Elekes--R\'onyai nonspecial and suppose its quadratic part is nonsingular:
\[
\boxed{4ac-b^2\neq0.}
\]

\begin{conjecture}[Nonsingular Quadratic Compression]
There exists \(c_f>0\) and finite sets
\[
A_j\subset\mathbb R,
\qquad
|A_j|\to\infty,
\]
such that
\[
\boxed{
|f(A_j,A_j)|\le |A_j|^{2-c_f}.
}
\]
\end{conjecture}

This would show that fixed power-saving compression is not merely a pathology caused by rank-one quadratic degeneration.

\section{Relation to the Unit-Distance and Sum-Product Constructions}

The common principle appears to be
\[
\boxed{
\text{obtain a fixed gain per algebraic coordinate}
+
\text{keep all losses exponential in the field degree}.
}
\]

Then, because
\[
d\asymp \log |A|,
\]
an exponential gain
\[
\theta^d,
\qquad
\theta<1,
\]
becomes a fixed power saving
\[
|A|^{-c}.
\]

In the unit-distance proof, the gain comes from many norm-one elements obtained from completely split primes in growing-degree fields.

In the recent sum-product construction, the gain comes from simultaneous additive and multiplicative compression in a high-dimensional number-field setting.

For the known polynomial counterexamples, the gain comes from local square-residue restrictions amplified through completely split primes.

For
\[
x^2+xy+y^2,
\]
the challenge is to find a \emph{new source of per-coordinate gain} compatible with its full-rank norm-form structure.

That is the main research question.

\section{Proposed Immediate Objective}

The immediate objective is therefore
\[
\boxed{
\text{Construct arbitrarily large }A\subset\mathbb R
\text{ such that }
|x^2+xy+y^2(A,A)|\le |A|^{2-c}.
}
\]

The first stage should not be to attack the full nonsingular quadratic conjecture, but to understand which pieces of the unit-distance and sum-product number-field constructions can survive for the Eisenstein norm.

In particular:

\begin{enumerate}[label=\arabic*.]
    \item Identify an arithmetic analogue of split-prime amplification that does not rely on rank-one degeneration.
    \item Test whether the factorization
    \[
    x^2+xy+y^2=N(x-y\omega)
    \]
    permits a construction using ideals, norm-one ratios, or correlated residue conditions.
    \item Quantify every gain and discriminant/container loss at the level of one field degree.
    \item Search for a fixed inequality of the form
    \[
    \text{arithmetic gain}>\text{lattice/discriminant cost}.
    \]
    \item Only after obtaining such a per-coordinate inequality, let the field degree tend to infinity to convert it into the desired fixed exponent \(c>0\).
\end{enumerate}

Thus the concrete conjecture \(x^2+xy+y^2\) is the first milestone, while the broader nonsingular quadratic statement provides the longer-term framework.

\end{document}